
\subsection{Die Dirac-Gleichung}
\label{subsec:diracgleichung}

Die im vorherigen Abschnitt hergeleitete Dirac-Gleichung%
\index{Dirac-Gleichung} beschreibt ein freies Elektron gleichzeitig
quantenmechanisch und relativistisch.

Sie lautet

\begin{equation}
	\mathrm{i}\hbar
	\frac{\partial\Psi}{\partial t}
	=
	\left(
	-\mathrm{i}\hbar c
	\boldsymbol{\alpha}\cdot\nabla
	+
	\beta m_{\mathrm e}c^2
	\right)
	\Psi.
	\label{eq:diracgleichung_hamiltonform_kapitel07}
\end{equation}

Diese Gleichung ähnelt auf den ersten Blick der zeitabhängigen
Schrödingergleichung:

\begin{equation}
	\mathrm{i}\hbar
	\frac{\partial\Psi}{\partial t}
	=
	\hat{H}\Psi.
\end{equation}

Der entscheidende Unterschied liegt im Hamiltonoperator. In der
Dirac-Gleichung lautet er

\begin{equation}
	\hat{H}_{\mathrm D}
	=
	-\mathrm{i}\hbar c
	\boldsymbol{\alpha}\cdot\nabla
	+
	\beta m_{\mathrm e}c^2.
	\label{eq:dirac_hamiltonoperator_kapitel07}
\end{equation}

Er enthält nicht nur den Impuls des Elektrons, sondern auch dessen
Ruheenergie und eine besondere Matrizenstruktur.

\begin{PhysicsBox}[Die Dirac-Gleichung]
	\label{box:diracgleichung_kapitel07}
	
	Die Dirac-Gleichung
	
	\begin{equation}
		\mathrm{i}\hbar
		\frac{\partial\Psi}{\partial t}
		=
		\left(
		-\mathrm{i}\hbar c
		\boldsymbol{\alpha}\cdot\nabla
		+
		\beta m_{\mathrm e}c^2
		\right)
		\Psi
	\end{equation}
	
	ist eine relativistische Bewegungsgleichung für ein Elektron.
	
	Sie verbindet
	
	\begin{itemize}
		\item die zeitliche Entwicklung eines Quantenzustandes,
		\item die Bewegung des Elektrons,
		\item seine Ruheenergie und
		\item seine Spinstruktur.
	\end{itemize}
\end{PhysicsBox}

\subsubsection*{Bedeutung der einzelnen Größen}

In der Dirac-Gleichung bezeichnet

\begin{itemize}
	\item \(\mathrm{i}\) die imaginäre Einheit,
	\item \(\hbar\) das reduzierte Plancksche Wirkungsquantum,
	\item \(c\) die Lichtgeschwindigkeit,
	\item \(m_{\mathrm e}\) die Ruhemasse des Elektrons,
	\item \(\nabla\) die räumliche Ableitung,
	\item \(\boldsymbol{\alpha}\) eine Gruppe von Matrizen,
	\item \(\beta\) eine weitere Matrix und
	\item \(\Psi\) den vollständigen Zustand des Elektrons.
\end{itemize}

Der Ausdruck

\begin{equation}
	-\mathrm{i}\hbar\nabla
\end{equation}

entspricht dem Impulsoperator%
\index{Impulsoperator}. Der erste Term des Hamiltonoperators beschreibt
daher die Bewegung des Elektrons.

Der zweite Term

\begin{equation}
	\beta m_{\mathrm e}c^2
\end{equation}

enthält die Ruheenergie

\begin{equation}
	E_0
	=
	m_{\mathrm e}c^2.
\end{equation}

Für das Elektron beträgt sie

\begin{equation}
	m_{\mathrm e}c^2
	\approx
	\SI{511}{\kilo\electronvolt}.
	\label{eq:ruheenergie_elektron_dirac}
\end{equation}

Die Matrix \(\beta\) ist notwendig, damit sowohl positive als auch
negative Energielösungen auftreten können.

\begin{DidacticBox}[Mehr als eine Bewegungsgleichung]
	\label{box:mehr_als_bewegungsgleichung_dirac}
	
	Die Dirac-Gleichung beschreibt nicht nur, wie sich ein Elektron durch
	den Raum bewegt.
	
	Ihre mathematische Struktur enthält zugleich
	
	\begin{itemize}
		\item die Ruheenergie,
		\item den Spin,
		\item das magnetische Moment und
		\item die Möglichkeit von Antiteilchen.
	\end{itemize}
	
	Diese Eigenschaften wurden nicht nachträglich angefügt. Sie ergeben
	sich aus der Gleichung selbst.
\end{DidacticBox}

\subsubsection*{Eine Wellenfunktion mit vier Komponenten}

In der einfachsten Form der Schrödingergleichung wird der Zustand eines
Teilchens durch eine einzelne komplexe Wellenfunktion beschrieben.

Bei der Dirac-Gleichung reicht eine einzelne Funktion nicht mehr aus.
Der Zustand des Elektrons besitzt vier Komponenten:

\begin{equation}
	\Psi
	=
	\begin{pmatrix}
		\psi_1 \\
		\psi_2 \\
		\psi_3 \\
		\psi_4
	\end{pmatrix}.
	\label{eq:vierkomponentiger_dirac_spinor_kapitel07}
\end{equation}

Eine solche Größe wird als Dirac-Spinor%
\index{Dirac-Spinor}\index{Spinor!vierkomponentiger} bezeichnet.

Die vier Komponenten sind keine vier getrennten Elektronen. Sie
gehören gemeinsam zu einem einzigen Quantenzustand.

Zwei Komponenten stehen mit den beiden möglichen Spinrichtungen des
Elektrons in Verbindung. Die vollständige vierkomponentige Struktur
wird zusätzlich benötigt, um positive und negative Energielösungen zu
beschreiben.

\begin{DidacticBox}[Vier Komponenten, ein Elektron]
	\label{box:vier_komponenten_ein_elektron}
	
	Der Dirac-Spinor besteht aus vier mathematischen Komponenten.
	
	Diese Komponenten beschreiben nicht vier Teilchen, sondern
	unterschiedliche Anteile eines einzigen relativistischen
	Elektronenzustandes.
	
	Sie enthalten Informationen über
	
	\begin{itemize}
		\item die Spinrichtung und
		\item die Teilchen-Antiteilchen-Struktur.
	\end{itemize}
\end{DidacticBox}

\subsubsection*{Die Rolle der Matrizen}

Die Größen \(\boldsymbol{\alpha}\) und \(\beta\) sind
Dirac-Matrizen\index{Dirac-Matrizen}. Sie wirken auf die vier
Komponenten des Spinors.

Gewöhnliche Zahlen würden lediglich jede Komponente mit demselben Wert
multiplizieren. Matrizen können dagegen die einzelnen Komponenten
miteinander koppeln.

Dadurch wird die zeitliche Entwicklung einer Komponente auch von den
anderen Komponenten beeinflusst.

Die Matrizen sind so konstruiert, dass beim zweimaligen Anwenden des
Dirac-Hamiltonoperators wieder die relativistische
Energie-Impuls-Beziehung entsteht:

\begin{equation}
	E^2
	=
	p^2c^2
	+
	m_{\mathrm e}^2c^4.
	\label{eq:energie_impuls_aus_diracgleichung}
\end{equation}

Damit ist sichergestellt, dass die Dirac-Gleichung mit der speziellen
Relativitätstheorie vereinbar ist.

\begin{MathBox}[Aufgabe der Dirac-Matrizen]
	\label{box:aufgabe_dirac_matrizen}
	
	Die Dirac-Matrizen erfüllen zwei Aufgaben:
	
	\begin{enumerate}
		\item Sie koppeln die vier Komponenten der Wellenfunktion.
		\item Sie sorgen dafür, dass die relativistische
		Energie-Impuls-Beziehung erfüllt wird.
	\end{enumerate}
	
	Die genaue Form und Algebra dieser Matrizen ist für das Verständnis der
	physikalischen Aussage zunächst nicht erforderlich.
\end{MathBox}

\subsubsection*{Der Spin folgt aus der Gleichung}

Der Spin\index{Spin!Dirac-Gleichung} musste in der
Schrödingergleichung zusätzlich eingeführt werden.

In der Dirac-Gleichung ist er dagegen bereits in der
vierkomponentigen Struktur und in den Matrizen enthalten.

Für das Elektron ergibt sich der Spin

\begin{equation}
	s
	=
	\frac{\hbar}{2}.
	\label{eq:spin_aus_diracgleichung}
\end{equation}

Bei einer Messung entlang einer gewählten Raumrichtung können zwei
mögliche Ergebnisse auftreten:

\begin{equation}
	s_z
	=
	+\frac{\hbar}{2}
\end{equation}

oder

\begin{equation}
	s_z
	=
	-\frac{\hbar}{2}.
\end{equation}

Diese beiden Möglichkeiten werden häufig vereinfacht als
Spin nach oben und Spin nach unten bezeichnet.

Dabei handelt es sich nicht um eine klassische Drehung des Elektrons.
Der Spin ist eine intrinsische Quanteneigenschaft.

\begin{PhysicsBox}[Der Spin ist keine Zusatzannahme]
	\label{box:spin_keine_zusatzannahme_dirac}
	
	In der Dirac-Theorie wird der Spin des Elektrons nicht nachträglich
	eingesetzt.
	
	Er folgt aus der relativistischen mathematischen Struktur der
	Gleichung.
	
	Das war einer der wichtigsten Erfolge der Dirac-Gleichung.
\end{PhysicsBox}

\subsubsection*{Das magnetische Moment}

Ein Teilchen mit elektrischer Ladung und Spin besitzt ein
magnetisches Moment%
\index{magnetisches Moment!Elektron}.

Das Elektron verhält sich dadurch in einem Magnetfeld ähnlich wie ein
kleiner Magnet. Seine möglichen Energiezustände hängen von der
Orientierung des Spins relativ zum Magnetfeld ab.

Für das magnetische Moment des Elektrons ergibt die Dirac-Theorie in
erster Näherung den Zusammenhang

\begin{equation}
	\boldsymbol{\mu}_{\mathrm s}
	=
	-g
	\frac{e}{2m_{\mathrm e}}
	\mathbf{S},
	\label{eq:magnetisches_moment_dirac}
\end{equation}

wobei \(g\) der sogenannte \(g\)-Faktor%
\index{g-Faktor!Elektron} ist.

Aus der Dirac-Gleichung folgt

\begin{equation}
	g
	=
	2.
	\label{eq:g_faktor_dirac}
\end{equation}

Experimente zeigen einen geringfügig abweichenden Wert. Diese kleine
Abweichung wird später durch die Quantenelektrodynamik erklärt.

\begin{DidacticBox}[Ein bemerkenswerter Erfolg]
	\label{box:g_faktor_erfolg_dirac}
	
	Die Dirac-Gleichung sagt für den \(g\)-Faktor des Elektrons den Wert
	
	\begin{equation}
		g=2
	\end{equation}
	
	voraus.
	
	Dieser Wert folgt aus der Theorie und muss nicht zusätzlich eingesetzt
	werden.
	
	Die geringe experimentelle Abweichung von \(2\) entsteht durch
	Wechselwirkungen des Elektrons mit dem quantisierten
	elektromagnetischen Feld.
\end{DidacticBox}

\subsubsection*{Positive und negative Energien}

Für ein freies Elektron liefert die relativistische
Energie-Impuls-Beziehung zwei mögliche Energiezweige:

\begin{equation}
	E
	=
	\pm
	\sqrt{
		p^2c^2
		+
		m_{\mathrm e}^2c^4
	}.
	\label{eq:positive_negative_energie_dirac}
\end{equation}

Der positive Zweig lautet

\begin{equation}
	E
	=
	+
	\sqrt{
		p^2c^2
		+
		m_{\mathrm e}^2c^4
	},
\end{equation}

der negative Zweig

\begin{equation}
	E
	=
	-
	\sqrt{
		p^2c^2
		+
		m_{\mathrm e}^2c^4
	}.
\end{equation}

Für ein ruhendes Teilchen mit \(p=0\) ergeben sich die Energien

\begin{equation}
	E
	=
	\pm m_{\mathrm e}c^2.
	\label{eq:ruheenergien_dirac}
\end{equation}

Die positiven Lösungen konnten unmittelbar mit Elektronen verbunden
werden. Die Bedeutung der negativen Lösungen war zunächst unklar.

Sie konnten jedoch nicht einfach gestrichen werden, weil sie
zwangsläufig aus der Gleichung folgen.

\begin{PhysicsBox}[Zwei Energiezweige]
	\label{box:zwei_energiezweige_dirac}
	
	Die Dirac-Gleichung besitzt positive und negative Energielösungen:
	
	\begin{equation}
		E
		=
		\pm
		\sqrt{
			p^2c^2
			+
			m_{\mathrm e}^2c^4
		}.
	\end{equation}
	
	Die negativen Lösungen waren kein Rechenfehler. Sie waren ein Hinweis
	darauf, dass die Theorie mehr als nur das Elektron beschreibt.
\end{PhysicsBox}

\subsubsection*{Die Wahrscheinlichkeitsdichte}

Eine Quantentheorie muss angeben, mit welcher Wahrscheinlichkeit ein
Teilchen an einem bestimmten Ort gefunden werden kann.

Für den Dirac-Spinor lautet die Wahrscheinlichkeitsdichte%
\index{Wahrscheinlichkeitsdichte!Dirac-Gleichung}

\begin{equation}
	\rho
	=
	\Psi^\dagger\Psi.
	\label{eq:wahrscheinlichkeitsdichte_dirac_kapitel07}
\end{equation}

Das Symbol \(\Psi^\dagger\) bezeichnet den komplex konjugierten und
transponierten Spinor.

Für die vier Komponenten ergibt sich

\begin{equation}
	\rho
	=
	\lvert\psi_1\rvert^2
	+
	\lvert\psi_2\rvert^2
	+
	\lvert\psi_3\rvert^2
	+
	\lvert\psi_4\rvert^2.
	\label{eq:wahrscheinlichkeitsdichte_dirac_komponenten}
\end{equation}

Da jeder einzelne Summand nichtnegativ ist, gilt

\begin{equation}
	\rho
	\geq
	0.
\end{equation}

Die Dirac-Gleichung besitzt damit eine sinnvoll interpretierbare
Wahrscheinlichkeitsdichte.

\begin{PhysicsBox}[Positive Wahrscheinlichkeitsdichte]
	\label{box:positive_wahrscheinlichkeitsdichte_dirac}
	
	Für die Dirac-Gleichung gilt
	
	\begin{equation}
		\rho
		=
		\Psi^\dagger\Psi.
	\end{equation}
	
	Diese Größe ist stets größer oder gleich null.
	
	Damit kann sie als Wahrscheinlichkeitsdichte für den Ort des Elektrons
	interpretiert werden.
\end{PhysicsBox}

\subsubsection*{Der nichtrelativistische Grenzfall}

Eine neue Theorie muss die ältere Theorie in dem Bereich wiedergeben,
in dem diese erfolgreich ist.

Für Geschwindigkeiten, die wesentlich kleiner als die
Lichtgeschwindigkeit sind,

\begin{equation}
	v
	\ll
	c,
\end{equation}

geht die Dirac-Gleichung näherungsweise in die
Pauli-Gleichung\index{Pauli-Gleichung!Grenzfall der Dirac-Gleichung}
über.

Werden zusätzlich Spin und Magnetfelder vernachlässigt, ergibt sich
schließlich wieder die Schrödingergleichung.

Der Zusammenhang kann vereinfacht dargestellt werden als

\begin{equation}
	\text{Dirac-Gleichung}
	\longrightarrow
	\text{Pauli-Gleichung}
	\longrightarrow
	\text{Schrödingergleichung}.
	\label{eq:grenzfall_dirac_pauli_schroedinger}
\end{equation}

\begin{DidacticBox}[Die ältere Theorie bleibt erhalten]
	\label{box:aeltere_theorie_bleibt_erhalten_dirac}
	
	Die Dirac-Gleichung ersetzt die Schrödingergleichung nicht in jedem
	praktischen Fall.
	
	Für langsame Elektronen liefert sie näherungsweise wieder die bekannten
	nichtrelativistischen Gleichungen.
	
	Die Schrödingergleichung bleibt deshalb für Atome, Moleküle und viele
	Festkörperprobleme weiterhin die zweckmäßigste Beschreibung.
\end{DidacticBox}

\subsubsection*{Was die Dirac-Gleichung erklärt}

Die Dirac-Gleichung brachte mehrere zuvor getrennte Eigenschaften des
Elektrons in eine gemeinsame Theorie.

Sie erklärt oder beschreibt

\begin{enumerate}
	\item die relativistische Bewegung des Elektrons,
	\item den Spin \(s=\hbar/2\),
	\item die beiden möglichen Spinrichtungen,
	\item das magnetische Moment,
	\item den \(g\)-Faktor \(g=2\) in erster Näherung,
	\item positive und negative Energielösungen und
	\item eine positive Wahrscheinlichkeitsdichte.
\end{enumerate}

\begin{HistoryBox}[Eine Gleichung mit unerwarteten Folgen]
	\label{box:unerwartete_folgen_diracgleichung}
	
	Paul Dirac\index{Dirac, Paul} wollte ursprünglich eine relativistische
	Quantengleichung für das Elektron finden.
	
	Die gefundene Gleichung leistete jedoch wesentlich mehr:
	
	\begin{itemize}
		\item Sie erklärte den Spin des Elektrons.
		\item Sie lieferte das magnetische Moment mit hoher Genauigkeit.
		\item Sie enthielt negative Energielösungen.
	\end{itemize}
	
	Gerade die zunächst problematischen negativen Energien führten zu einer
	der bedeutendsten Vorhersagen der modernen Physik: der Existenz des
	Positrons.
\end{HistoryBox}

Die Dirac-Gleichung beschreibt das Elektron daher nicht nur genauer als
die Schrödingergleichung. Sie zeigt, dass Spin, Relativität und
Antiteilchen eng miteinander verbunden sind.

Die physikalische Bedeutung der negativen Energielösungen und die
Vorhersage des Positrons werden im folgenden Abschnitt genauer
betrachtet.
